\begin{question}
  Consider a three-dimensional ket space.  If a certain set of
  orthonormal kets, say, $\ket{1}$, $\ket{2}$, $\ket{3}$, are used as
  the base kets, the operators $A$ and $B$ are represented by

  \begin{equation*}
    \begin{alignedat}{2}
      &A &\dot{=}
      \begin{pmatrix}
        a & 0 & 0\\
        0 & -a & 0\\
        0 & 0 & -a\\
      \end{pmatrix}
      , &
      \quad
      & B& \dot{=}
      \begin{pmatrix}
        b & 0 & 0\\
        0 & 0 & -ib\\
        0 & ib & 0\\
      \end{pmatrix}
    \end{alignedat}
  \end{equation*}

  with $a$ and $b$ both real.

  \begin{questionpart}
    Obviously $A$ exhibits a degenerate spectrum.  Does $B$ also
    exhibit a degenerate spectrum?
  \end{questionpart}

  \begin{questionpart}
    Show that $A$ and $B$ commute.
  \end{questionpart}

  \begin{questionpart}
    Find a new set of orthonormal kets which are simultaneous
    eigenkets of both $A$ and $B$.  Specify the eigenvalues of $A$ and
    $B$ for each of the three eigenkets.  Does your specification of
    eigenvalues completely characterize each eigenket?
  \end{questionpart}
\end{question}

\begin{purpose}
  Should probably note the anti-diagonal portion of $B$.  We've seen
  that pattern before with the Pauli spin operators.  The main purpose
  of this problem seems to be to get some practice with finding
  eigenvectors in the presence of degeneracy.

  Warning to the reader: I'm hitting some pretty serious senioritis
  here after having finished problem \ref{1.24:theQuestion} which was
  the capstone problem for the chapter so... yeah... be a bit patient
  with me please.
\end{purpose}

\begin{answer}

  \begin{answerpart}{More degenerates?}

    It's worth noting that $B^{\dagger} = B \implies \lambda \in {R}$
    and the same is true for $A$.  That means that $\lambda$'s must be
    real.

    \begin{equation}
      \begin{alignedat}{2}
        &&0 &= \text{Det}(B - \lambda I)\\
        &&&=
        \text{Det}
        \begin{pmatrix}
          b - \lambda & 0 & 0\\
          0 & -\lambda & -ib\\
          0 & ib & -\lambda\\
        \end{pmatrix}
        \\
        &&&=b\
        \text{Det}
        \begin{pmatrix}
          1 - \lambda' & 0 & 0\\
          0 & -\lambda' & -i\\
          0 & i & -\lambda'\\
        \end{pmatrix}
        \\
        &&&= b(1 - \lambda')(-\lambda')(-\lambda') - (ib)(-ib)(b - \lambda')\\
        &&&= b(1 - \lambda')(\lambda'^2 - 1^2)\\
        &&&= (b - \lambda)(\lambda^2 - b^2)\\
      \end{alignedat}
    \end{equation}

    We can easily see that the eigenvalues are $\lambda = \pm b$ with
    two of them being $b$, ladies and germs, we have ourselves some
    degeneracy.  When you apply these eigenvalues you'll want to have
    care to normalize things since $B$ has a factor of $b$ in there.
  \end{answerpart}

  \begin{answerpart}{Morning Commute}

    \begin{equation}
      \begin{split}
        AB
        &=
        \begin{pmatrix}
          a & 0 & 0\\
          0 & -a & 0\\
          0 & 0 & -a\\
        \end{pmatrix}
        \begin{pmatrix}
          b & 0 & 0\\
          0 & 0 & -ib\\
          0 & ib & 0\\
        \end{pmatrix}
        \\
        &= \begin{pmatrix}
          ab & 0 & 0\\
          0 & 0 & iab\\
          0 & -iab & 0\\
        \end{pmatrix}
        \\
        &= ab\begin{pmatrix}
          1 & 0 & 0\\
          0 & 0 & i\\
          0 & -i & 0\\
        \end{pmatrix}
        \\
      \end{split}
    \end{equation}
    \begin{equation}
      \begin{split}
        BA
        &=
        \begin{pmatrix}
          b & 0 & 0\\
          0 & 0 & -ib\\
          0 & ib & 0\\
        \end{pmatrix}
        \begin{pmatrix}
          a & 0 & 0\\
          0 & -a & 0\\
          0 & 0 & -a\\
        \end{pmatrix}
        \\
        &= \begin{pmatrix}
          ab & 0 & 0\\
          0 & 0 & iab\\
          0 & -iab & 0\\
        \end{pmatrix}
        \\
        &= ab\begin{pmatrix}
          1 & 0 & 0\\
          0 & 0 & i\\
          0 & -i & 0\\
        \end{pmatrix}
        \\
      \end{split}
    \end{equation}

    \begin{equation}
      AB = BA \Rightarrow [A, B] = 0\qed
    \end{equation}

    So, if they commute, then there must exist a simultaneous
    eigenbasis.  We showed that earlier in
    problem...umm... \ref{1.17:theQuestion}.
  \end{answerpart}

  \begin{answerpart}{Perfect Union}

    Since we know the two operators commute, the easiest thing to
    articulate is reference equation \ref{1.17:result} in problem
    \ref{1.17:theQuestion} and rewrite the product of the operators
    as:

    \begin{equation}
      AB
      = \sum_i \lambda_{a,i}\lambda_{b,i}\Lambda_i
      = \sum_i \lambda_i\Lambda_i
    \end{equation}

    So if we just view thew product of the eigenvalues for the
    individual operators as a single eigenvalue for the product
    ($\lambda_i \equiv \lambda_{a,i}\lambda_{b,i}$) then we can find a
    simultaneous eigenbasis by finding the eigenbasis of the product
    of the operators.  Did you follow that?  No...that's probably
    because I'm only half paying attention to this one.

    \begin{equation}
      \begin{alignedat}{2}
        &&0
        &=\text{Det}
        \left(
        ab
        \begin{pmatrix}
          1-\lambda' & 0 & 0\\
          0 & -\lambda' & i\\
          0 & -i & -\lambda'\\
        \end{pmatrix}
        \right)
        \\
        \implies &&0 &= \text{Det}
        \begin{pmatrix}
          1-\lambda' & 0 & 0\\
          0 & -\lambda' & i\\
          0 & -i & -\lambda'\\
        \end{pmatrix}
        \\
        &&&= (1-\lambda')(-\lambda')(-\lambda') - (-i)(i)(1-\lambda')\\
        &&&= (1-\lambda')(\lambda'^2-1^2)\\
        &&&= (ab-\lambda')(\lambda'^2-(ab)^2)\\
      \end{alignedat}
    \end{equation}

    So, we can clearly see that the simultaneous eigenvalues are
    exactly the same as before (with scaling factors of $a$, $b$, and
    $ab$ for matrices $A$, $B$, and $AB/BA$ respectively).  With our
    eigenvalues of $\pm 1$, we can solve for the non-degenerate
    eigenvectors:

    \begin{equation}
      \begin{alignedat}{2}
        &&
        \begin{pmatrix}
          1 & 0 & 0\\
          0 & 0 & i\\
          0 & -i & 0\\
        \end{pmatrix}
        \begin{pmatrix}
          x\\ y\\ z\\
        \end{pmatrix}
        &= \begin{pmatrix}
          x\\ y\\ z\\
        \end{pmatrix}
        \\
        \implies
        &&x &= x\\
        &&iz &= y\\
        &&-iy &= z\\
      \end{alignedat}
    \end{equation}

    So, let's pick some values.  We can pick anything we want for $x$,
    so let's go with $0$.  Then we'll pick $1-i$ for $y$.  Then we'll
    also negate and that leaves us with the \emph{non-normalized}
    \emph{non-degenerate} eigenkets

    \begin{equation}
      \ket{1} \dot{=}
      \begin{bmatrix}
        0\\ 0\\ 0\\
      \end{bmatrix}
      ,\quad
      \ket{2} \dot{=}
      \frac{1}{2}
      \begin{bmatrix}
        0\\ 1+i\\ 1-i\\
      \end{bmatrix}
      ,\quad
      \ket{3} \dot{=}
      \frac{1}{2}
      \begin{bmatrix}
        0\\ 1-i\\ 1+i\\
      \end{bmatrix}
    \end{equation}

    OK, the problem is what to do about that degenerate eigenvalue.
    The problem asks for \emph{orthogonal} eigenvectors.  It's easy
    enough to manually construct one that works (like by taking a
    cross product), but according to my nagging assistant I should
    probably show the more methodical approach.

    Since we identified that the $x$ value there can be anything we
    want for the $\pm b$ eigenvalues, then we set it to 0 for those
    eigenvectors.  For our final eigenvector, we'll just choose a $1$
    in that spot.  It's also just easier that way so that makes our
    final set

    \begin{equation}
      \ket{1} \dot{=}
      \begin{bmatrix}
        1\\ 0\\ 0\\
      \end{bmatrix}
      ,\quad
      \ket{2} \dot{=}
      \frac{1}{2}
      \begin{bmatrix}
        0\\ 1+i\\ 1-i\\
      \end{bmatrix}
      ,\quad
      \ket{3} \dot{=}
      \frac{1}{2}
      \begin{bmatrix}
        0\\ 1-i\\ 1+i\\
      \end{bmatrix}
    \end{equation}

    Explicitly stating our eigenvalues for the various vectors

    \begin{equation}
      \begin{split}
        \lambda_{A,1} &= a\\
        \lambda_{A,2} &= -a\\
        \lambda_{A,3} &= -a\\
        \lambda_{B,1} &= b\\
        \lambda_{B,2} &= -b\\
        \lambda_{B,3} &= b\\
        \lambda_{AB,1} &= ab\\
        \lambda_{AB,2} &= ab\\
        \lambda_{AB,3} &= -ab\\
      \end{split}
    \end{equation}

    Remember to mind your manners and your normalization constants
    when you use these.  As for the final part of the question: did
    the eigenvalues fully specify the eigenvectors the answer they're
    looking for isn't about the mechanics of finding Eigenvectors,
    it's about the quantum mechanics of observing the observable.  If
    you observe the observable $B$ and get $b$ in return, it tells you
    that you're either in $\ket{1}$ or $\ket{3}$.  To figure out the
    state of the system, you have to simultaneously measure both $A$
    and $B$ where you get

    \begin{equation}
      \begin{split}
        (\lambda_{A,1}, \lambda_{B,1}) &= (a, b)\\
        (\lambda_{A,2}, \lambda_{B,2}) &= (-a, -b)\\
        (\lambda_{A,3}, \lambda_{B,3}) &= (-a, b)\\
      \end{split}
    \end{equation}

    Note that you can just as easily pick $-1$ for $\ket{1}$ which
    gets you exactly the same effective result.  The arbitrary nature
    of that selection is why you don't see the 4th theoretical state
    of $(-a, -b)$ above.  In the case of an observable, that's
    essentially a selection of ``Which way is up?'' or ``Are we going
    to use a left-handed or right-handed coordinate system?'' or
    ``Should we use emacs or vim?''.  The answers are ``Up is away
    from the enemy's gate'', ``right-handed'', and ``If you're a vim
    user then you are dead to me.'' respectively.

  \end{answerpart}

\end{answer}
